\slajdid{09_1-s015}
\begin{frame}[label=09_1-s015]
\gusttekst\setlength{\abovedisplayskip}{4pt}\setlength{\belowdisplayskip}{4pt}
\begin{columns}[c,onlytextwidth]\column{.25\textwidth}\centering
\begin{tikzpicture}[x=.70cm,y=.72cm,font=\sitneoznake,>=Latex,remember picture]\draw[->](-.15,0)--(3.3,0)node[right,yshift=4pt]{$V_I$};\draw[->](0,-.1)--(0,2.7)node[above]{$V_O$};\node[left]at(0,2.35){$V_{DD}$};\node[below]at(2.8,0){$V_{DD}$};\draw(0,2.35)--(.35,2.35)..controls(.65,2.35)and(.7,2.2)..(.85,1.75)..controls(1.05,1.2)and(1.25,.3)..(1.55,.25)--(2.8,.15);\draw[dashed](2.8,-.1)--(2.8,2.65);\draw[dashed](.3,-.1)--(.3,2.65);\draw[dashed,DEcrvena](1.5,-.1)--(1.5,2.65);\node at(.22,2.6){1.};\node at(.87,2.6){2.};\node at(2.1,2.6){3.};\end{tikzpicture}

\column{.73\textwidth}
Kada ulazni napon još poraste tranzistor ulazi u triodnu oblast i tada je
\[\frac{V_{DD}-V_O}{R_L}=\frac{k_n}{2}\frac1{1+\frac{V_O}{L_nE_{Cn}}}\bigl(2V_O(V_I-V_{Tn})-V_O^2\bigr)\]
\end{columns}\medskip
\begin{columns}[c,onlytextwidth]\column{.25\textwidth}\centering
\begin{tikzpicture}[x=.70cm,y=.72cm,font=\sitneoznake,>=Latex,remember picture]\draw[->](-.15,0)--(3.3,0)node[right,yshift=4pt]{$V_I$};\draw[->](0,-.1)--(0,2.7)node[above]{$V_O$};\node[left]at(0,2.35){$V_{DD}$};\node[below]at(2.8,0){$V_{DD}$};\draw(0,2.35)--(.35,2.35)..controls(.65,2.35)and(.7,2.2)..(.85,1.75)..controls(1.05,1.2)and(1.25,.3)..(1.55,.25)--(2.8,.15);\draw[dashed](2.8,-.1)--(2.8,2.65);\draw[dashed,DEplava](.58,-.1)--(.58,2.65);\draw[dashed,DEplava](1.55,-.1)--(1.55,2.65);\coordinate(s15high)at(.58,2.31);\coordinate(s15low)at(1.55,.25);\node at(.3,2.6){1.};\node at(1.05,2.6){2.};\node at(2.15,2.6){3.};\end{tikzpicture}

\column{.73\textwidth}
1. Oblast – zakočen i u zasićenju\par
2. Oblast – u zasićenju i u triodnoj\par
3. Oblast – u triodnoj\par\medskip
\tikz[remember picture,baseline=(s15f1.base)]{\node[inner sep=0pt](s15f1){$\displaystyle\frac{V_{DD}-V_O}{R_L}=\frac{k_n}{2}\frac1{1+\frac{(V_I-V_{Tn})}{L_nE_{Cn}}}(V_I-V_{Tn})^2$};}\quad Pojačanje može biti manje od 1\par\medskip
\tikz[remember picture,baseline=(s15f2.base)]{\node[inner sep=0pt](s15f2){$\displaystyle\frac{V_{DD}-V_O}{R_L}=\frac{k_n}{2}\frac1{1+\frac{V_O}{L_nE_{Cn}}}\bigl(2V_O(V_I-V_{Tn})-V_O^2\bigr)$};}\par
\hfill Pojačanje može biti veće od 1
\tikz[remember picture,overlay]{\draw[->,DEplava](s15f1.west)--(s15high);\draw[->,DEplava](s15f2.west)--(s15low);}
\end{columns}
\end{frame}
